\documentclass[11pt,a4paper]{article}
\usepackage[margin=2.2cm]{geometry}
\usepackage[utf8]{inputenc}\usepackage[T1]{fontenc}
\usepackage{amsmath,amssymb}
\usepackage{fancyhdr}\usepackage{enumitem}\usepackage{xcolor}
\usepackage{hyperref}
\pagestyle{fancy}\fancyhf{}
\rhead{Math 1A Worksheet 4}
\lhead{Songpengfei}
\rfoot{Page \thepage}
\newcommand{\solvertag}[1]{{\footnotesize\color{gray}\textsf{[#1]}}}

\title{Homework Solutions}
\author{Songpengfei}
\date{2026-10-09}
\begin{document}
\maketitle

\noindent\textbf{Solved: 9/10 (90.0\%)}\quad\textit{Auto-generated by C4C homework-solver.}

\section*{Problem Q1}
Suppose $f$ is a function and $L$ and $a$ are real numbers.

\paragraph{Solution.} \solvertag{sympy\_limit} 

\begin{enumerate}[leftmargin=*]
  \item See sub-problems for individual limit computations.
\end{enumerate}

\paragraph{(a)} $\forall \varepsilon > 0, \exists \delta > 0: 0 < |x-a| < \delta \Rightarrow |f(x)-L| < \varepsilon$

\paragraph{(b)} $\text{See graph}$

\paragraph{(c)} $\text{See graph}$

\paragraph{(d)} $\text{See graph}$

\noindent\textbf{Answer:}\quad $\text{See sub-problems}$

\medskip

\section*{Problem Q2}
Is $\infty$ a number? What does $\lim_{x \to a} f(x) = \infty$ mean? Does $\lim_{x \to a} f(x) = \infty$ mean that the limit exists?

\paragraph{Solution.} \solvertag{conceptual\_template} 

\begin{enumerate}[leftmargin=*]
  \item $\infty$ is \textbackslash\{\}textbf\{not\} a real number.
  \item $\lim_{x \to a} f(x) = \infty$ means that $f(x)$ grows without bound as $x \to a$.
  \item The limit does \textbackslash\{\}textbf\{not exist\} in the usual sense; we say the limit ``equals infinity'' as a shorthand for divergence.
\end{enumerate}

\noindent\textbf{Answer:}\quad $\infty \text{ is not a number. The limit does not exist.}$

\medskip

\section*{Problem Q3}
What does the Squeeze Theorem say? Draw a picture illustrating this theorem.

\paragraph{Solution.} \solvertag{conceptual\_template} 

\begin{enumerate}[leftmargin=*]
  \item \textbackslash\{\}textbf\{Squeeze Theorem:\} If $g(x) \le f(x) \le h(x)$ for all $x$ near $a$ (except possibly at $a$), and $\lim_{x \to a} g(x) = \lim_{x \to a} h(x) = L$, then $\lim_{x \to a} f(x) = L$.
  \item Intuitively: if $f$ is ``squeezed'' between two functions that both converge to $L$, then $f$ must also converge to $L$.
\end{enumerate}

\noindent\textbf{Answer:}\quad $\text{If } g(x) \le f(x) \le h(x) \text{ and } \lim g = \lim h = L, \text{ then } \lim f = L.$

\medskip

\section*{Problem Q4}
What do we mean by $\lim_{x \to a^-} f(x) = L$? What is meant by $\lim_{x \to a^+} f(x) = L$? Graph a function $y = f(x)$ and some point $x = a$ where $\lim_{x \to a^-} f(x) \neq \lim_{x \to a^+} f(x)$. Find an algebraic expression for such a function.

\paragraph{Solution.} \solvertag{conceptual\_template} 

\begin{enumerate}[leftmargin=*]
  \item $\lim_{x \to a^-} f(x) = L$: the \textbackslash\{\}textbf\{left-hand limit\}. As $x$ approaches $a$ from values \textbackslash\{\}emph\{less than\} $a$, $f(x)$ approaches $L$.
  \item $\lim_{x \to a^+} f(x) = L$: the \textbackslash\{\}textbf\{right-hand limit\}. As $x$ approaches $a$ from values \textbackslash\{\}emph\{greater than\} $a$, $f(x)$ approaches $L$.
  \item The two-sided limit $\lim_{x \to a} f(x) = L$ exists if and only if both one-sided limits exist and are equal.
  \item Example: $f(x) = \begin{cases} x & x < 0 \\ x+1 & x \ge 0 \end{cases}$ at $a = 0$: $\lim_{x \to 0^-} f(x) = 0 \neq 1 = \lim_{x \to 0^+} f(x)$.
\end{enumerate}

\noindent\textbf{Answer:}\quad $\lim_{x \to a^-} f(x): \text{from below}; \quad \lim_{x \to a^+} f(x): \text{from above}$

\medskip

\section*{Problem P1}
(a) Is the limit of a sum always the sum of the limits? Give an example where $\lim_{x \to a}[f(x) + g(x)]$ exists even though neither $\lim_{x \to a} f(x)$ nor $\lim_{x \to a} g(x)$ exist.

\paragraph{Solution.} \solvertag{conceptual\_template} 

\begin{enumerate}[leftmargin=*]
  \item No, the limit of a sum is \textbackslash\{\}emph\{not\} always the sum of the limits.
  \item The \textbackslash\{\}textbf\{Sum Law\} requires both individual limits to exist: $\lim[f(x)+g(x)] = \lim f(x) + \lim g(x)$ \textbackslash\{\}emph\{if both limits exist\}.
  \item \textbackslash\{\}textbf\{Counterexample:\} Let $f(x) = \frac{1}{x}$ and $g(x) = -\frac{1}{x}$. Neither $\lim_{x \to 0} f(x)$ nor $\lim_{x \to 0} g(x)$ exists, but $\lim_{x \to 0}[f(x)+g(x)] = \lim_{x \to 0} 0 = 0$.
\end{enumerate}

\noindent\textbf{Answer:}\quad $\text{No. } f(x)=1/x,\; g(x)=-1/x: \lim[f+g]=0 \text{ but } \lim f, \lim g \text{ DNE}$

\medskip

\section*{Problem P2}
(a) Find $\lim_{x \to 0^+}(\frac{1}{x} - \frac{1}{|x|})$, $\lim_{x \to 0^-}(\frac{1}{x} - \frac{1}{|x|})$, and $\lim_{x \to 0}(\frac{1}{x} - \frac{1}{|x|})$, if they exist.

\paragraph{Solution.} \solvertag{sympy} 

\begin{enumerate}[leftmargin=*]
  \item $\lim_{x \to 0^+} - \frac{1}{\left|{x}\right|} + \frac{1}{x} = 0$
  \item $\lim_{x \to 0^-} - \frac{1}{\left|{x}\right|} + \frac{1}{x} = -\infty$
  \item $\lim_{x \to 0} - \frac{1}{\left|{x}\right|} + \frac{1}{x} = 0$
\end{enumerate}

\paragraph{(b)} $0$

\paragraph{(c)} $\text{DNE}$

\noindent\textbf{Answer:}\quad $0$

\medskip

\section*{Problem P3}
For what kinds of functions can one evaluate the limit $\lim_{x \to a} f(x)$ by just plugging in $x = a$? Give an example of a function $f(x)$ for which the limit $\lim_{x \to a} f(x)$ exists and can be evaluated, yet for which this substitution doesn't work.

\paragraph{Solution.} \solvertag{conceptual\_template} 

\begin{enumerate}[leftmargin=*]
  \item \textbackslash\{\}textbf\{Continuous functions\} allow direct substitution: $\lim_{x \to a} f(x) = f(a)$.
  \item This includes polynomials, rational functions (at non-zero denominator points), trigonometric, exponential, and logarithmic functions at points in their domain.
  \item Example where limit exists but substitution fails: $f(x) = \frac{x^2 - 1}{x - 1}$ at $x = 1$. Direct substitution gives $\frac{0}{0}$, but $\lim_{x \to 1} \frac{x^2-1}{x-1} = \lim_{x \to 1}(x+1) = 2$.
\end{enumerate}

\noindent\textbf{Answer:}\quad $\text{Continuous functions: } \lim_{x \to a} f(x) = f(a)$

\medskip

\section*{Problem AP1}
Below are listed three limits. Match the best method to each limit and evaluate each limit.

\paragraph{Solution.} \solvertag{sympy\_limit} 

\begin{enumerate}[leftmargin=*]
  \item See sub-problems for individual limit computations.
\end{enumerate}

\paragraph{(a)} $-20$

\paragraph{(b)} $10$

\paragraph{(c)} $0$

\noindent\textbf{Answer:}\quad $\text{See sub-problems}$

\medskip

\section*{Problem AP2}
Using a calculator, try to guess $\lim_{n \to \infty} n \sin(\frac{\pi}{n})$. (Remember that $n$ is in radians, not degrees!)

\paragraph{Solution.} \solvertag{sympy} 

\begin{enumerate}[leftmargin=*]
  \item $\lim_{x \to \infty} n \sin{\left(\frac{\pi}{n} \right)} = \pi$
\end{enumerate}

\noindent\textbf{Answer:}\quad $\pi$

\medskip

\section*{Problem AP3}
Draw a circle of radius 1. Inscribe a regular $n$-gon inside it.

{\color{red}\textit{Not solved.} 确定性求解器未覆盖，且本次运行禁用了 LLM（--no-llm）}

\medskip

\end{document}
